\scanpage{178}
\mspage{91}
% L01–02
(3) de GH1 :
\[
\forall s,t\in S,\ \text{on a}\quad\{s,t\}\in A_0\Longleftrightarrow\{f^{-1}s,f^{-1}\lambda t\}\in A_0. \tag*{(58)}
\]

% L03–09
\deleted{et donc en} Cette condition ne change pas si on \deleted{\unsure{rempl}} remplace $f$ par $fg$ \deleted{\unsure{$fg$}}, avec $g\in G_0$, donc elle ne dépend en effet que de l'image $\dot f$ de $f$ dans $G/G_0\simeq Z'$ $(\overset{\sim}{\leftarrow} Z)$. Quitte à changer $f$ mod $G_0$, \unsure{SPG|SPS} $f\in B'$. Donc, la vertu essentielle de l'isom $\alpha:Z=\widetilde B''/\widetilde U\longrightarrow Z'=B'/\widetilde U$ qui ``corrige'' $\kappa$, c'est la validité de (58), pour \deleted{\unsure{$f$}} $\lambda\in Z$ et $f\in B'$ reliés par $\dot f=\alpha(\lambda)$.

% L10–12
\[
\begin{aligned}
&\forall\lambda\in Z,f\in B',s,t\in S,\ \text{tels que }\dot f=\alpha(\lambda)\ \text{dans }Z'=B'/\widetilde U,\\
&\text{on a l'équivalence}\\
&\qquad \{s,t\}\in A_0\Longleftrightarrow\{f^{-1}s,f^{-1}\lambda t\}\in A_0
\end{aligned}\tag*{(59)}
\]

% L13–18
Nous pouvons dans cette équivalence écrire
\[
s=us_0,\ t=vs_0\qquad\text{avec }u,v\in G_0,\quad\deleted{\text{et les deux}}.
\]
et les deux termes de l'équivalence dans (59) s'écrivent respectivement
\[
\left\{\begin{aligned}
\text{a)}\ &\exists w\in G_0,\quad w(s_0,t_0)=(us_0,vs_0)\\
\text{b)}\ &\exists w'\in G_0\quad w'(s_0,t_0)= f^{-1}us_0,\underbrace{f^{-1}\lambda vs_0}_{f^{-1}v\lambda s_0}\bigr),
\end{aligned}\right. \tag*{(60)}
\]
\margin{Ici $t_0$ est le deuxième sommet du repère d'épinglage fixé $(s_0,t_0,u_0)$}

% L19–24
soit encore, en introduisant $\rho\in G_0$ tel que l'on ait
\[
t_0=\rho s_0\qquad(\rho\in G_0) \tag*{(61)}
\]
ces termes s'écrivent aussi respectivement
\[
\left\{\begin{aligned}
\text{a)}\ &\exists w\in G_0,\quad w^{-1}u\in\widetilde U,\ (w\rho)^{-1}v\in\widetilde U\\
\text{b)}\ &\exists w'\in G_0,\quad w'^{-1}f^{-1}u\in B',\ (w'\rho)^{-1}f^{-1}\lambda v\in B'
\end{aligned}\right. \tag*{(62)}
\]
Donc la condition à vérifier s'écrit

% archive sheet number: 177
